\begin{textsample}{POS Dim 1 – 2010 – Score 29.00 – 2010/t002739.txt}  \label{ex:f1_pos_027}
I ' ve been fascinated for a lifetime by \textbf{the} beauty, form and function of giant bluefin tuna.
Bluefin are warmblooded like us.
They ' re \textbf{the} largest of \textbf{the} tunas, \textbf{the} second-largest \textbf{fish} in \textbf{the} \textbf{sea} — bony \textbf{fish}.
They actually are a \textbf{fish} that is endothermic — powers through \textbf{the} ocean with warm muscles like a mammal.
That ' s one of our bluefin at \textbf{the} Monterey Bay Aquarium.
You can see in its shape and its streamlined design it ' s powered for ocean swimming.
It flies through \textbf{the} ocean on its pectoral \textbf{fins}, gets lift, powers its movements with a lunate tail.
It ' s actually got a naked skin for most of its body, so it reduces friction with \textbf{the} water.
This is what one of nature ' s finest machines.
Now, bluefin were revered by Man for all of human history.
For 4, 000 years, we \textbf{fished} sustainably for this animal, and it ' s evidenced in \textbf{the} art that we see from thousands of years ago.
Bluefin are in cave paintings in France.
They ' re on coins that date back 3, 000 years.
This \textbf{fish} was revered by humankind.
It was \textbf{fished} sustainably till all of time, except for our generation.
Bluefin are pursued wherever they go — there is a gold rush on Earth, and this is a gold rush for bluefin.
There are traps that \textbf{fish} sustainably up until recently.
And yet, \textbf{the} type of \textbf{fishing} going on today, with pens, with enormous stakes, is really \textbf{wiping} bluefin ecologically off \textbf{the} \textbf{planet}.
Now bluefin, in general, goes to one place : Japan.
Some of you may be guilty of having contributed to \textbf{the} demise of bluefin.
They ' re delectable muscle, rich in \textbf{fat} — absolutely taste delicious.
And that ' s their problem ; we ' re eating them to death.
Now in \textbf{the} Atlantic, \textbf{the} story is pretty simple.
Bluefin have two populations : one large, one small. \textbf{The} North American population is \textbf{fished} at about 2, 000 ton. \textbf{The} European population and North African — \textbf{the} \textbf{Eastern} bluefin tuna — is \textbf{fished} at tremendous levels : 50, 000 tons over \textbf{the} last decade almost every year. \textbf{The} result is whether you ' re looking at \textbf{the} West or \textbf{the} Eastern bluefin population, there ' s been tremendous decline on both sides, as much as 90 percent if you go back with your baseline to 1950.
For that, bluefin have been given a status equivalent to \textbf{tigers}, to \textbf{lions}, to certain African \textbf{elephants} and to pandas.
These \textbf{fish} have been proposed for an endangered species listing in \textbf{the} past two months.
They were voted on and rejected just two weeks ago, despite outstanding science that shows from two committees this \textbf{fish} meets \textbf{the} criteria of CITES I.
And if it ' s tunas you don ' t care about, perhaps you might be interested that international long lines and pursing chase down tunas and bycatch animals such as leatherbacks, sharks, marlin, albatross.
These animals and their demise \textbf{occurs} in \textbf{the} tuna fisheries. \textbf{The} challenge we face is that we know very little about tuna, and everyone in \textbf{the} room knows what it looks like when an African \textbf{lion} takes down its prey.
I doubt anyone has seen a giant bluefin feed.
This tuna symbolizes what ' s \textbf{the} problem for all of us in \textbf{the} room.
It ' s \textbf{the} 21st century, but we really have only just begun to really study our oceans in a deep way.
Technology has come of age that ' s allowing us to see \textbf{the} Earth from space and go deep into \textbf{the} \textbf{seas} remotely.
And we ' ve got to use these technologies immediately to get a better understanding of how our ocean realm works.
Most of us from \textbf{the} ship — even I — look out at \textbf{the} ocean and see this homogeneous \textbf{sea}.
We don ' t know where \textbf{the} structure is.
We can ' t tell where are \textbf{the} watering holes like we can on an African plain.
We can ' t see \textbf{the} corridors, and we can ' t see what it is that brings together a tuna, a leatherback and an albatross.
We ' re only just beginning to understand how \textbf{the} physical oceanography and \textbf{the} biological oceanography come together to create a seasonal force that actually causes \textbf{the} upwelling that might make a hot spot a hope spot. \textbf{The} reasons these challenges are great is that technically it ' s difficult to go to \textbf{sea}.
It ' s hard to study a bluefin on its turf, \textbf{the} entire Pacific realm.
It ' s really tough to get up close and personal with a mako shark and try to put a \textbf{tag} on it.
And then imagine being Bruce Mate ' s team from OSU, getting up close to a \textbf{blue} \textbf{whale} and fixing a \textbf{tag} on \textbf{the} \textbf{blue} \textbf{whale} that stays, an engineering challenge we ' ve yet to really overcome.
So \textbf{the} story of our team, a dedicated team, is \textbf{fish} and chips.
We basically are taking \textbf{the} same satellite phone parts, or \textbf{the} same parts that are in your \textbf{computer}, chips.
We ' re putting them together in unusual ways, and this is taking us into \textbf{the} ocean realm like never before.
And for \textbf{the} first time, we ' re able to watch \textbf{the} journey of a tuna beneath \textbf{the} ocean using light and \textbf{photons} to measure sunrise and sunset.
Now, I ' ve been working with tunas for over 15 years.
I have \textbf{the} privilege of being a partner with \textbf{the} Monterey Bay Aquarium.
We ' ve actually taken a sliver of \textbf{the} ocean, put it behind glass, and we together have put bluefin tuna and yellowfin tuna on display.
When \textbf{the} veil of bubbles lifts every morning, we can actually see a community from \textbf{the} Pelagic ocean, one of \textbf{the} only places on Earth you can see giant bluefin swim by.
We can see in their beauty of form and function, their ceaseless activity.
They ' re flying through their space, ocean space.
And we can bring two million people a year into contact with this \textbf{fish} and show them its beauty.
Behind \textbf{the} scenes is a working lab at Stanford University partnered with \textbf{the} Monterey Bay Aquarium.
Here, for over 14 or 15 years, we ' ve actually brought in both bluefin and yellowfin in captivity.
We ' d been studying these \textbf{fish}, but first we had to learn how to husbandry them.
What do they like to eat?
What is it that they ' re happy with?
We go in \textbf{the} tanks with \textbf{the} tuna — we touch their naked skin — it ' s pretty amazing.
It feels wonderful.
And then, better yet, we ' ve got our own version of tuna whisperers, our own Chuck Farwell, Alex Norton, who can take a big tuna and in one motion, put it into an envelope of water, so that we can actually work with \textbf{the} tuna and learn \textbf{the} \textbf{techniques} it takes to not injure this \textbf{fish} who never sees a boundary in \textbf{the} open \textbf{sea}.
Jeff and Jason there, are scientists who are going to take a tuna and put it in \textbf{the} equivalent of a treadmill, a flume.
And that tuna thinks it ' s going to Japan, but it ' s staying in place.
We ' re actually measuring its oxygen \textbf{consumption}, its energy \textbf{consumption}.
We ' re taking this data and building better models.
And when I see that tuna — this is my favorite view — I begin to wonder : how did this \textbf{fish} solve \textbf{the} longitude problem before we did?
So take a look at that animal.
That ' s \textbf{the} closest you ' ll probably ever get.
Now, \textbf{the} activities from \textbf{the} lab have taught us now how to go out in \textbf{the} open ocean.
So in a program called Tag-A-Giant we ' ve actually gone from Ireland to Canada, from Corsica to Spain.
We ' ve \textbf{fished} with many nations around \textbf{the} world in an effort to basically put electronic \textbf{computers} inside giant tunas.
We ' ve actually \textbf{tagged} 1, 100 tunas.
And I ' m going to show you three \textbf{clips}, because I \textbf{tagged} 1, 100 tunas.
It ' s a very hard process, but it ' s a ballet.
We bring \textbf{the} tuna out, we measure it.
A team of fishers, captains, scientists and technicians work together to keep this animal out of \textbf{the} ocean for about four to five minutes.
We put water over its gills, give it oxygen.
And then with a lot of effort, after \textbf{tagging}, putting in \textbf{the} \textbf{computer}, making sure \textbf{the} stalk is sticking out so it senses \textbf{the} environment, we send this \textbf{fish} back into \textbf{the} \textbf{sea}.
And when it goes, we ' re always happy.
We see a flick of \textbf{the} tail.
And from our data that gets collected, when that \textbf{tag} comes back, because a fisher returns it for a thousand-dollar reward, we can get tracks beneath \textbf{the} \textbf{sea} for up to five years now, on a backboned animal.
Now sometimes \textbf{the} tunas are really large, such as this \textbf{fish} off Nantucket.
But that ' s about half \textbf{the} size of \textbf{the} biggest tuna we ' ve ever \textbf{tagged}.
It takes a human effort, a team effort, to bring \textbf{the} \textbf{fish} in.
In this case, what we ' re going to do is put a pop-up satellite archival \textbf{tag} on \textbf{the} tuna.
This \textbf{tag} rides on \textbf{the} tuna, senses \textbf{the} environment around \textbf{the} tuna and actually will come off \textbf{the} \textbf{fish}, detach, float to \textbf{the} \textbf{surface} and send back to Earth-orbiting satellites position data estimated by math on \textbf{the} \textbf{tag}, pressure data and temperature data.
And so what we get then from \textbf{the} pop-up satellite \textbf{tag} is we get away from having to have a human interaction to recapture \textbf{the} \textbf{tag}.
Both \textbf{the} electronic \textbf{tags} I ' m talking about are expensive.
These \textbf{tags} have been engineered by a variety of teams in North America.
They are some of our finest instruments, our new technology in \textbf{the} ocean today.
One community in general has given more to help us than any other community.
And that ' s \textbf{the} fisheries off \textbf{the} state of North Carolina.
There are two villages, Harris and Morehead City, every winter for over a decade, held a party called Tag-A-Giant, and together, fishers worked with us to \textbf{tag} 800 to 900 \textbf{fish}.
In this case, we ' re actually going to measure \textbf{the} \textbf{fish}.
We ' re going to do something that in recent years we ' ve started : take a mucus sample.
Watch how shiny \textbf{the} skin is ; you can see my reflection there.
And from that mucus, we can get gene profiles, we can get information on gender, checking \textbf{the} pop-up \textbf{tag} one more time, and then it ' s out in \textbf{the} ocean.
And this is my favorite.
With \textbf{the} help of my former postdoc, Gareth Lawson, this is a gorgeous picture of a single tuna.
This tuna is actually moving on a numerical ocean. \textbf{The} warm is \textbf{the} Gulf Stream, \textbf{the} cold up there in \textbf{the} Gulf of Maine.
That ' s where \textbf{the} tuna wants to go — it wants to forage on schools of herring — but it can ' t get there.
It ' s too cold.
But then it warms up, and \textbf{the} tuna pops in, gets some \textbf{fish}, maybe comes back to home base, goes in again and then comes back to winter down there in North Carolina and then on to \textbf{the} Bahamas.
And my favorite scene, three tunas going into \textbf{the} Gulf of Mexico.
Three tunas \textbf{tagged}.
Astronomically, we ' re \textbf{calculating} positions.
They ' re coming together.
That could be tuna sex — and there it is.
That is where \textbf{the} tuna spawn.
So from data like this, we ' re able now to put \textbf{the} map up, and in this map you see thousands of positions generated by this decade and a half of \textbf{tagging}.
And now we ' re showing that tunas on \textbf{the} western side go to \textbf{the} \textbf{eastern} side.
So two populations of tunas — that is, we have a Gulf population, one that we can \textbf{tag} — they go to \textbf{the} Gulf of Mexico, I showed you that — and a second population.
Living amongst our tunas — our North American tunas — are European tunas that go back to \textbf{the} Med.
On \textbf{the} hot spots — \textbf{the} hope spots — they ' re mixed populations.
And so what we ' ve done with \textbf{the} science is we ' re showing \textbf{the} International Commission, building new models, showing them that a two-stock no-mixing model — to this day, used to reject \textbf{the} CITES treaty — that model isn ' t \textbf{the} right model.
This model, a model of overlap, is \textbf{the} way to move forward.
So we can then predict where management places should be.
Places like \textbf{the} Gulf of Mexico and \textbf{the} Mediterranean are places where \textbf{the} single species, \textbf{the} single population, can be captured.
These become forthright in places we need to protect. \textbf{The} center of \textbf{the} Atlantic where \textbf{the} mixing is, I could imagine a policy that lets Canada and America \textbf{fish}, because they manage their fisheries well, they ' re doing a good job.
But in \textbf{the} international realm, where \textbf{fishing} and overfishing has really gone wild, these are \textbf{the} places that we have to make hope spots in.
That ' s \textbf{the} size they have to be to protect \textbf{the} bluefin tuna.
Now in a second project called \textbf{Tagging} of Pacific Pelagics, we took on \textbf{the} \textbf{planet} as a team, those of us in \textbf{the} Census of Marine Life.
And, funded primarily through Sloan Foundation and others, we were able to actually go in, in our project — we ' re one of 17 field programs and begin to take on \textbf{tagging} large numbers of predators, not just tunas.
So what we ' ve done is actually gone up to \textbf{tag} salmon shark in Alaska, met salmon shark on their home territory, followed them catching salmon and then went in and figured out that, if we take a salmon and put it on a line, we can actually take up a salmon shark — This is \textbf{the} cousin of \textbf{the} white shark — and very carefully — note, I say"very carefully,"— we can actually keep it calm, put a hose in its mouth, keep it off \textbf{the} deck and then \textbf{tag} it with a satellite \textbf{tag}.
That satellite \textbf{tag} will now have your shark phone home and send in a message.
And that shark leaping there, if you look carefully, has an antenna.
It ' s a free swimming shark with a satellite \textbf{tag} jumping after salmon, sending home its data.
Salmon sharks aren ' t \textbf{the} only sharks we \textbf{tag}.
But there goes salmon sharks with this meter-level resolution on an ocean of temperature — warm colors are warmer.
Salmon sharks go down to \textbf{the} tropics to pup and come into Monterey.
Now right next door in Monterey and up at \textbf{the} Farallones are a white shark team led by Scott Anderson — there — and Sal Jorgensen.
They can throw out a target — it ' s a carpet shaped like a \textbf{seal} — and in will come a white shark, a curious critter that will come right up to our 16-ft. boat.
It ' s a several thousand-pound animal.
And we ' ll wind in \textbf{the} target.
And we ' ll place an acoustic \textbf{tag} that says,"OMSHARK 10165,"or something like that, acoustically with a ping.
And then we ' ll put on a satellite \textbf{tag} that will give us \textbf{the} long- distance journeys with \textbf{the} light-based geolocation algorithms solved on \textbf{the} \textbf{computer} that ' s on \textbf{the} \textbf{fish}.
So in this case, Sal ' s looking at two \textbf{tags} there, and there they are : \textbf{the} white sharks of California going off to \textbf{the} white shark cafe and coming back.
We also \textbf{tag} makos with our NOAA colleagues, \textbf{blue} sharks.
And now, together, what we can see on this ocean of color that ' s temperature, we can see ten-day worms of makos and salmon sharks.
We have white sharks and \textbf{blue} sharks.
For \textbf{the} first time, an ecoscape as large as ocean-scale, showing where \textbf{the} sharks go. \textbf{The} tuna team from TOPP has done \textbf{the} unthinkable : three teams \textbf{tagged} 1, 700 tunas, bluefin, yellowfin and albacore all at \textbf{the} same time — carefully rehearsed \textbf{tagging} programs in which we go out, pick up juvenile tunas, put in \textbf{the} \textbf{tags} that actually have \textbf{the} sensors, stick out \textbf{the} tuna and then let them go.
They get returned, and when they get returned, here on a NASA numerical ocean you can see bluefin in \textbf{blue} go across their corridor, returning to \textbf{the} Western Pacific.
Our team from UCSC has \textbf{tagged} \textbf{elephant} \textbf{seals} with \textbf{tags} that are glued on their heads, that come off when they slough.
These \textbf{elephant} \textbf{seals} cover half an ocean, take data down to 1, 800 feet — amazing data.
And then there ' s Scott Shaffer and our shearwaters wearing tuna \textbf{tags}, light-based \textbf{tags}, that now are going to take you from New Zealand to Monterey and back, journeys of 35, 000 nautical miles we had never seen before.
But now with light-based geolocation \textbf{tags} that are very small, we can actually see these journeys.
Same thing with Laysan albatross who travel an entire ocean on a trip sometimes, up to \textbf{the} same zone \textbf{the} tunas use.
You can see why they might be caught.
Then there ' s George Schillinger and our leatherback team out of Playa Grande \textbf{tagging} leatherbacks that go right past where we are.
And Scott Benson ' s team that showed that leatherbacks go from Indonesia all \textbf{the} way to Monterey.
So what we can see on this moving ocean is we can finally see where \textbf{the} predators are.
We can actually see how they ' re using ecospaces as large as an ocean.
And from this information, we can begin to map \textbf{the} hope spots.
So this is just three years of data right here — and there ' s a decade of this data.
We see \textbf{the} \textbf{pulse} and \textbf{the} seasonal activities that these animals are going on.
So what we ' re able to do with this information is boil it down to hot spots, 4, 000 deployments, a huge herculean task, 2, 000 \textbf{tags} in an area, shown here for \textbf{the} first time, off \textbf{the} California coast, that appears to be a gathering place.
And then for sort of an encore from these animals, they ' re helping us.
They ' re carrying instruments that are actually taking data down to 2, 000 meters.
They ' re taking information from our \textbf{planet} at very critical places like Antarctica and \textbf{the} Poles.
Those are \textbf{seals} from many countries being released who are sampling underneath \textbf{the} ice sheets and giving us temperature data of oceanographic quality on both poles.
This data, when visualized, is captivating to watch.
We still haven ' t figured out best how to visualize \textbf{the} data.
And then, as these animals swim and give us \textbf{the} information that ' s important to climate issues, we also think it ' s critical to get this information to \textbf{the} public, to engage \textbf{the} public with this kind of data.
We did this with \textbf{the} Great Turtle Race — \textbf{tagged} \textbf{turtles}, brought in four million hits.
And now with Google ' s Oceans, we can actually put a white shark in that ocean.
And when we do and it swims, we see this magnificent bathymetry that \textbf{the} shark knows is there on its path as it goes from California to Hawaii.
But maybe Mission Blue can fill in that ocean that we can ' t see.
We ' ve got \textbf{the} capacity, NASA has \textbf{the} ocean.
We just need to put it together.
So in conclusion, we know where Yellowstone is for North America ; it ' s off our coast.
We have \textbf{the} technology that ' s shown us where it is.
What we need to think about perhaps for Mission Blue is increasing \textbf{the} biologging capacity.
How is it that we can actually take this type of activity elsewhere?
And then finally — to basically get \textbf{the} message home — maybe use live links from animals such as \textbf{blue} \textbf{whales} and white sharks.
Make killer apps, if you will.
A lot of people are excited when sharks actually went under \textbf{the} Golden Gate Bridge.
Let ' s connect \textbf{the} public to this activity right on their iPhone.
That way we do away with a few internet myths.
So we can save \textbf{the} bluefin tuna.
We can save \textbf{the} white shark.
We have \textbf{the} science and technology.
Hope is here.
Yes we can.
We need just to apply this capacity further in \textbf{the} oceans.
Thank you.

% matched lemmas: blue, calculate, clip, computer, consumption, eastern, elephant, fat, fin, fish, fishing, lion, occur, photon, planet, pulse, sea, seal, surface, tag, technique, the, tiger, turtle, whale, wipe
\end{textsample}
